\begin{theorem}[有界纤维满射的最小字母自然扩张实现]\label{thm:xi-natural-extension-minimal-d-alphabet-realization}
设 \(T:X\twoheadrightarrow X\) 为可数集合上的有限纤维满射，并记
\[
D:=\sup_{y\in X}\#T^{-1}(y)<\infty,
\qquad
[D]:=\{1,\dots,D\}.
\]
对每个 \(y\in X\) 任选单射
\[
\iota_y:T^{-1}(y)\hookrightarrow [D],
\qquad
\beta(x):=\iota_{T(x)}(x).
\]
定义编码
\[
C_D:X^{\leftarrow}\longrightarrow X\times [D]^{\NN},
\]
\[
C_D\bigl((x_0,x_{-1},x_{-2},\dots)\bigr)
:=
\bigl(x_0,(\beta(x_{-1}),\beta(x_{-2}),\dots)\bigr),
\]
并记
\[
\widehat X_D:=C_D(X^{\leftarrow})\subset X\times [D]^{\NN}.
\]
则 \(C_D\) 把自然扩张 \((X^{\leftarrow},\sigma)\) 共轭到 \(\widehat X_D\) 上的双射
\[
\widehat T_D\bigl(x,(a_1,a_2,\dots)\bigr)
:=
\bigl(Tx,(\beta(x),a_1,a_2,\dots)\bigr).
\]
若进一步存在某个 \(y_0\in X\) 满足
\[
\#T^{-1}(y_0)=D,
\]
则任何此类右移插入实现都不可能把字母表缩小到大小 \(D-1\) 的集合。因而 \(D\) 正是自然扩张单寄存器实现的最小可行字母数。
\end{theorem}
\begin{proof}
由定理 \ref{thm:killo-natural-extension-branch-register}，若允许目标字母表为 \(\NN\)，则上述编码把自然扩张 \((X^{\leftarrow},\sigma)\) 共轭为右移插入系统。由于每条纤维的大小都不超过 \(D\)，全部单射 \(\iota_y\) 可统一取值于 \([D]\)，故像空间自动落在 \(X\times [D]^{\NN}\) 中，从而得到所示子空间 \(\widehat X_D\) 与双射 \(\widehat T_D\)。

最小性则正是命题 \ref{prop:xi-natural-extension-single-register-minimal-alphabet} 的内容：若某条纤维大小等于 \(D\)，则任何单寄存器实现的字母表大小都至少为 \(D\)；反之取字母表 \([D]\) 即可实现。因此最小可行字母数恰为 \(D\)。证毕。
\end{proof}

\endinput
